\begin{result}{A.4.40}{Exact local sizes and arithmetic type of the invariant leaf}{Proved}
Let $E/K$ be an elliptic curve in the algebraic normalization of
A.4.14, let $a\in K$, and write $d=[K:\Q]$.  In its universal vector
extension $G=E^\dagger$, let $\widehat C_a$ be the formal leaf at the
identity of the invariant line $L_a=\{s=au\}$.  Choose, once over $K$,
\'etale parameters $X,Y$ at the identity such that
\[
 X(e)=Y(e)=0,\qquad D_aX(e)=1,\qquad D_aY(e)=0.
 \tag{A.4.40a}
\]
Write the actual formal graph as
$Y=\phi_a(X)=\sum_{n\ge2}b_nX^n\in K[[X]]$.

Outside a finite set of places, put
\[
 \mathcal B_a=\{v:K_{p_v}+aH_{p_v}\ne0\pmod v\},\qquad
 d_v=[K_v:\Q_{p_v}].
\]
For the absolute value normalized by $\|p\|_v=p^{-d_v}$, the geometric
size of this precise embedded formal curve satisfies
\[
 S_v(\widehat C_a)=\|p\|_v^{1/(p-1)}\quad(v\in\mathcal B_a),
 \qquad
 S_v(\widehat C_a)\ge\|p\|_v^{1/(p(p-1))}
       \quad(v\notin\mathcal B_a).
 \tag{A.4.40b}
\]
Consequently its adelic size product is zero.  In particular this
formal leaf fails the A-analyticity condition in Bost--Chambert-Loir's
Definition~3.7.

For Andr\'e's normalization $|p|_v=p^{-d_v/d}$ and his arithmetic
coefficient invariant $\tau$, the same fixed graph satisfies
\[
 \tau(\phi_a)\ \ge\
 \limsup_{T\to\infty}\frac1T
       \sum_{\substack{v\in\mathcal B_a\\p_v\le T}}
                       \frac{d_v}{d}\log p_v
 \ \ge\ \frac1{4d}>0.
 \tag{A.4.40c}
\]
Thus the zero-$\tau$ hypothesis of Andr\'e's Theorem~5.4.3 fails for
this actual graph.  These statements hold for every algebraic $a$,
including the established CM Hodge coefficients.  They do not decide
whether the actual non-CM coefficient $h$ can be algebraic.
\end{result}

\begin{proof}
We use the actual integral formal embedding and track its coefficients
at the first degree divisible by the residue characteristic.

\paragraph{The fixed coordinate map.}
An invertible $K$-linear change of any \'etale parameters gives
(A.4.40a), since $D_a(e)\ne0$.  Spread the smooth group, its invariant
vector field, the identity, and these parameters over
$\mathcal O_K[1/N]$, increasing $N$ so that the vector field is
primitive and the parameter map is \'etale at the identity.
This is one fixed model and coordinate choice; it does not vary with
$v$.  Its map to $\mathbb A^2$ identifies the formal leaf with the
specified graph.  The geometric sizes agree by
Bost--Chambert-Loir's Theorem~3.1(b), printed p.~10.
Their size is also unchanged by an integral linear coordinate
automorphism.\cite{AuditBostChambertLoir2008v3}

\paragraph{The coefficient of degree $p$.}
Exclude the finitely many ramified primes of $K$ and those dividing
$N$, and let $p=p_v>3$.  In characteristic zero the formal flow of
$D_a$ gives
\[
 X(t)=\sum_{n\ge1}\frac{D_a^nX(e)}{n!}t^n,
 \qquad
 Y(t)=\sum_{n\ge1}\frac{D_a^nY(e)}{n!}t^n.
 \tag{A.4.40d}
\]
All coefficients below degree $p$ are integral, $X'(0)=1$, and
$Y'(0)=0$.  The compositional inverse $g$ of $X(t)$ has integral
coefficients below degree $p$.  In $Y(g(X))$, the possibly
nonintegral degree-$p$ coefficient of $g$ cannot enter through the
linear term, because that term of $Y$ is zero.  All contributions
from degrees $2,\ldots,p-1$ use only coefficients of $g$ below
degree $p$.  The degree-$p$ term of $Y$ contributes unchanged.
Therefore
\[
 b_n\in\mathcal O_v\quad(n<p),\qquad
 b_p-\frac{D_a^pY(e)}{p!}\in\mathcal O_v.
 \tag{A.4.40e}
\]

Reduction of the invariant vector field commutes with its $p$-fold
iterate.  In characteristic $p$ that iterate is an invariant
derivation.  Hence failure of the invariant line to be $p$-closed is
equivalent to its transverse component being nonzero at the
identity.  By (A.4.40a), this means that $D_a^pY(e)$ is a unit.
A.4.31 identifies these places exactly with $\mathcal B_a$.
Since $p!$ has precisely one factor $p$,
\[
                  \|b_p\|_v=\|p\|_v^{-1}
                         \qquad(v\in\mathcal B_a).
 \tag{A.4.40f}
\]
No assertion about the denominators of other degrees is needed.

\paragraph{From the coefficient to geometric size.}
The integral shear $(X,Y)\mapsto(X,Y+X)$ has determinant one and
changes the graph to $X+\phi_a(X)$, with unit linear coefficient
and unchanged higher coefficients.  Local analytic existence for
the nonsingular differential equation gives a positive convergence
radius.  Bost--Chambert-Loir's Proposition~3.5(2), printed
pp.~11--12, identifies its size with the minimum of $1$ and the
exponential of the first Newton slope.  Its degree-$p$ coefficient
therefore gives the upper bound
\[
           S_v\le\|b_p\|_v^{-1/(p-1)}
                 =\|p\|_v^{1/(p-1)}.
\]
Their Proposition~3.6, printed p.~12, gives the reverse inequality
for every integral nonsingular rank-one foliation.  At an absolutely
unramified $p$-closed place the same proposition gives the improved
lower bound in (A.4.40b).  All its hypotheses hold on the chosen
identity neighborhood.\cite{AuditBostChambertLoir2008v3}

\paragraph{The adelic sum.}
The bounds just proved imply
\[
 \sum_{p_v\le T}\log S_v^{-1}
   =\sum_{\substack{v\in\mathcal B_a\\p_v\le T}}
                  d_v\frac{\log p_v}{p_v-1}+O(1).
 \tag{A.4.40g}
\]
Here the nonnegative contribution of the other good places is
bounded by $d\sum_p\log p/(p(p-1))$, a convergent series.  The
discarded finite set changes only the bounded term.  A.4.32 states
that the complementary weighted upper density is at least
$1/(4d)$.  Since the total weighted sum is asymptotic to
$d\log T$, the right side of (A.4.40g) is unbounded.  Thus the
product is zero, which is exactly the failure of condition~(ii) in
the source's Definition~3.7, printed p.~13.  The model independence
and \'etale invariance in that definition and Theorem~3.1 ensure
that this conclusion concerns $\widehat C_a\subset G$, not only
its coordinate presentation.\cite{AuditBostChambertLoir2008v3}

\paragraph{Andr\'e's coefficient invariant.}
Put $\epsilon_v=d_v/d$.  In the normalization of Andr\'e's
\S5.1, printed p.~70, his \S5.3, printed p.~74, defines
\[
 \tau(\phi_a)=\inf_{P_0}\limsup_{n\to\infty}\frac1n
    \sum_{p_v\ge P_0}\max_{m\le n}\log^+|b_m|_v.
 \tag{A.4.40h}
\]
The order of the infimum, limit superior and sum is part of this
formula.\cite{AuditAndre2004Dwork}
At a place $v\in\mathcal B_a$ with $p_v\le n$, (A.4.40f) contributes
$\epsilon_v\log p_v$ to that sum.  Set
\[
 \theta_{\mathcal B}(T)=
     \sum_{\substack{v\in\mathcal B_a\\p_v\le T}}
                        \epsilon_v\log p_v,
 \qquad
 H_{\mathcal B}(T)=
     \sum_{\substack{v\in\mathcal B_a\\p_v\le T}}
                  \epsilon_v\frac{\log p_v}{p_v-1}.
\]
A fixed prime threshold removes a bounded amount from
$\theta_{\mathcal B}(n)$, so (A.4.40h) gives
$\tau(\phi_a)\ge\limsup\theta_{\mathcal B}(T)/T$.
This limit superior is precisely Andr\'e's weighted upper natural
density by his Lemma~5.1.1, printed pp.~70--71.

For completeness, its comparison with the density used in A.4.32
has a direct proof.  Replacing $p-1$ by $p$ changes
$H_{\mathcal B}$ by a bounded amount.  Partial summation gives
\[
 H_{\mathcal B}(T)=\frac{\theta_{\mathcal B}(T)}T
            +\int_2^T\frac{\theta_{\mathcal B}(t)}{t^2}\,dt+O(1).
\]
If $r=\limsup\theta_{\mathcal B}(T)/T$, then for every
$\varepsilon>0$ the integrand is at most
$(r+\varepsilon)/t$ beyond a fixed bound.  Hence
\[
 \limsup\frac{H_{\mathcal B}(T)}{\log T}\le r.
\]
The left side is at least $1/(4d)$ by A.4.32, because
$\sum_{v\mid p}\epsilon_v=1$.  This proves (A.4.40c) using only
upper densities; it supplies no lower-density assertion for
$\mathcal B_a$.
\end{proof}

\paragraph{Source scope and reproduction.}
The Bost--Chambert-Loir source is the opened 55-page arXiv version~3
of 2~October~2008.  The Andr\'e source is the original 2004 publisher
preview containing all the cited printed pages, not a claim to
have obtained the entire chapter.  Both agents visually checked
printed pp.~70, 74 and~76, including the product over places in
Theorem~5.4.3 and the formula (A.4.40h).
That theorem requires $\tau=0$, finite $\rho$, and simultaneous
meromorphic uniformizations with product of radii greater than one.
The established positive $\tau$ already excludes this direct
substitution; no unproved equality with a connection invariant is
used.  Likewise Theorem~6.1 of Bost--Chambert-Loir requires
A-analyticity, which (A.4.40g) explicitly disproves for these leaves.
The actual analytic uniformization
$u\mapsto[(u,au)]$ makes the archimedean canonical tangent seminorm
zero by their Proposition~5.15, printed p.~33, applied to discs of
arbitrarily large radius; it does not supply the failed finite-place
product.

The exact certificate
\texttt{check\_leaf\_prime\_graph\_coefficient.py}
works in the quotient by $t^{p+1}$ for $p=5,7$.  It constructs and
checks the inverse series and verifies that the degree-$p$
coefficient minus $y_p/p$ is an integer polynomial in the lower
coefficients, independent of $x_p$.  The general argument is
(A.4.40e); the two finite certificates reproduce its decisive
coefficient step.  The code, retained output and metadata identify
the exact inputs, hashes and runtime.
